% questions.tex —— 2009 年·普通高等学校招生考试·山东卷·文科 · 题目内容片段（无答案）
% 由 tools/generate_projects.py 依来源总账自 2009/2009_数学真题_山东文_试题.md 生成，勿手改（改源 MD 后重跑）


状态：needs-review；依固定候选 PDF 两页（印刷页 547–548）逐题核验，图形已矢量重绘，正式构建与逐页验收见回执。源未标分值，不补分值。第 17 题原页条件「在 $x=\pi$ 处取最小值」与所给函数不相容（见题注与解析），按原页照录并登记问题。

\gkpart{一、选择题}

\begin{question}
\gkti{1}集合 $A=\{0,2,a\}$，$B=\{1,a^2\}$．若 $A\cup B=\{0,1,2,4,16\}$，则 $a$ 的值为（　）

（A）$0$

（B）$1$

（C）$2$

（D）$4$

\end{question}
\begin{question}
\gkti{2}复数 $\dfrac{3-\mathrm i}{1-\mathrm i}$ 等于（　）

（A）$1+2\mathrm i$

（B）$1-2\mathrm i$

（C）$2+\mathrm i$

（D）$2-\mathrm i$

\end{question}
\begin{question}
\gkti{3}将函数 $y=\sin2x$ 的图象向左平移 $\dfrac{\pi}{4}$ 个单位，再向上平移 $1$ 个单位，所得图象的函数解析式是（　）

（A）$y=\cos2x$

（B）$y=2\cos^2x$

（C）$y=1+\sin\left(2x+\dfrac{\pi}{4}\right)$

（D）$y=2\sin^2x$

\Needspace{0.4\textheight}
\end{question}
\begin{question}
\gkti{4}一空间几何体的三视图如图所示，则该几何体的体积为（　）

（A）$2\pi+2\sqrt3$

（B）$4\pi+2\sqrt3$

（C）$2\pi+\dfrac{2\sqrt3}{3}$

（D）$4\pi+\dfrac{2\sqrt3}{3}$
\nopagebreak[4]
\begin{center}% 2009 年山东卷文科第 4 题：三视图（圆柱与其上四棱锥）。
% 来源 PDF：普通高考/2009/2009山东文.pdf，SHA256 1807901a21da57e3d45382294189b0c8e795722b6879842319ae75428639fad6
% 原页：第 1 页（A3 三栏）第 1 栏中部，印刷页 547。
% 重绘依据：正视图为 2×2 矩形上方接顶角三角形（斜棱标 2），侧视图同形；
%   俯视图为半径 1 的圆内接正方形（对角线为两条互相垂直的直径）；尺寸按原图标注。
\begin{tikzpicture}[scale=0.8, semithick, every node/.style={inner sep=1.3pt, font=\small}]
  % 正视图
  \begin{scope}[shift={(0,0)}]
    \draw (0,0) rectangle (2,2);
    \draw (0,2) -- (1,3.0) -- (2,2);
    \node[anchor=east] at (-0.12,1) {$2$};
    \node[anchor=north] at (1,-0.12) {$2$};
    \node[anchor=south east] at (0.42,2.55) {$2$};
    \node[anchor=north] at (1,-0.42) {正（主）视图};
  \end{scope}
  % 侧视图
  \begin{scope}[shift={(4.2,0)}]
    \draw (0,0) rectangle (2,2);
    \draw (0,2) -- (1,3.0) -- (2,2);
    \node[anchor=north] at (1,-0.12) {$2$};
    \node[anchor=south east] at (0.42,2.55) {$2$};
    \node[anchor=north] at (1,-0.42) {侧（左）视图};
  \end{scope}
  % 俯视图
  \begin{scope}[shift={(1.0,-4.6)}]
    \draw (1,1) circle (1);
    \draw (0,1) -- (2,1) -- (1,2) -- (1,0) -- cycle;
    \node[anchor=north] at (1,-0.18) {俯视图};
  \end{scope}
\end{tikzpicture}
\end{center}
\nopagebreak[4]
\end{question}
\begin{question}
\gkti{5}在 $\mathbf R$ 上定义运算 $\odot$：$a\odot b=ab+2a+b$，则满足 $x\odot(x-2)<0$ 的实数 $x$ 的取值范围为（　）

（A）$(0,2)$

（B）$(-2,1)$

（C）$(-\infty,-2)\cup(1,+\infty)$

（D）$(-1,2)$

\Needspace{0.4\textheight}
\end{question}
\begin{question}
\gkti{6}函数 $y=\dfrac{\mathrm e^x+\mathrm e^{-x}}{\mathrm e^x-\mathrm e^{-x}}$ 的图象大致为（　）
\nopagebreak[4]
\begin{center}% 2009 年山东卷文科第 6 题：y=(e^x+e^{-x})/(e^x-e^{-x}) 的四个候选图象。
% 来源 PDF：普通高考/2009/2009山东文.pdf，SHA256 1807901a21da57e3d45382294189b0c8e795722b6879842319ae75428639fad6
% 原页：第 1 页（A3 三栏）第 2 栏上部，印刷页 547。
% 重绘依据：(A) 双曲余切：右支由 +∞ 递减趋于 1，左支由 −1 递减到 −∞，虚线为 y=±1；
%   (B) 在 y 轴处有跳跃（(0,1) 与 (0,−1) 为空心点），整体递增；
%   (C) 过原点、以 y 轴为竖直切线、两侧递增（形如 y=x^{1/3}）；
%   (D) 关于 y 轴对称、在原点取 0、两侧递增（形如 y=|x|^{1/3}）。
\begin{tikzpicture}[semithick, every node/.style={inner sep=1.1pt, font=\small}]
  % (A)
  \begin{scope}[shift={(0,0)}, x=0.62cm, y=0.62cm]
    \draw[dashed] (-2.0,1) -- (2.4,1);
    \draw[dashed] (-2.0,-1) -- (2.4,-1);
    \draw[-{Stealth}] (-2.1,0) -- (2.5,0) node[right] {$x$};
    \draw[-{Stealth}] (0,-2.2) -- (0,2.2) node[above] {$y$};
    \begin{scope}
      \clip (-2.1,-2.2) rectangle (2.5,2.2);
      \draw[domain=0.14:2.3,samples=140,smooth,variable=\x]
        plot (\x, {(exp(\x)+exp(-\x))/(exp(\x)-exp(-\x))});
      \draw[domain=-2.3:-0.14,samples=140,smooth,variable=\x]
        plot (\x, {(exp(\x)+exp(-\x))/(exp(\x)-exp(-\x))});
    \end{scope}
    \node[anchor=north] at (0.02,-0.10) {$O$};
    \node[anchor=south east] at (-0.06,1) {$1$};
    \node[anchor=north east] at (-0.06,-1) {$-1$};
    \node at (0.6,-2.85) {(A)};
  \end{scope}
  % (B)
  \begin{scope}[shift={(5.0,0)}, x=0.62cm, y=0.62cm]
    \draw[-{Stealth}] (-2.1,0) -- (2.5,0) node[right] {$x$};
    \draw[-{Stealth}] (0,-2.2) -- (0,2.2) node[above] {$y$};
    \begin{scope}
      \clip (-2.1,-2.2) rectangle (2.5,2.2);
      \draw[domain=-2.2:-0.001,samples=100,smooth,variable=\x]
        plot (\x, {1+0.55*\x+0.10*\x*\x});
      \draw[domain=0.001:2.2,samples=100,smooth,variable=\x]
        plot (\x, {-1+0.55*\x+0.10*\x*\x});
    \end{scope}
    \draw[fill=white] (0,1) circle (1.4pt);
    \draw[fill=white] (0,-1) circle (1.4pt);
    \node[anchor=north] at (0.02,-0.10) {$O$};
    \node[anchor=west] at (0.08,1) {$1$};
    \node at (0.6,-2.85) {(B)};
  \end{scope}
  % (C)
  \begin{scope}[shift={(0,-5.4)}, x=0.62cm, y=0.62cm]
    \draw[-{Stealth}] (-2.1,0) -- (2.5,0) node[right] {$x$};
    \draw[-{Stealth}] (0,-2.2) -- (0,2.2) node[above] {$y$};
    \begin{scope}
      \clip (-2.1,-2.2) rectangle (2.5,2.2);
      \draw[domain=-2.2:2.2,samples=160,smooth,variable=\x]
        plot (\x, {sign(\x)*pow(abs(\x),1/3)*1.25});
    \end{scope}
    \node[anchor=north] at (0.02,-0.10) {$O$};
    \node[anchor=south west] at (1,-0.02) {$1$};
    \node at (0.6,-2.85) {(C)};
  \end{scope}
  % (D)
  \begin{scope}[shift={(5.0,-5.4)}, x=0.62cm, y=0.62cm]
    \draw[-{Stealth}] (-2.1,0) -- (2.5,0) node[right] {$x$};
    \draw[-{Stealth}] (0,-2.2) -- (0,2.2) node[above] {$y$};
    \begin{scope}
      \clip (-2.1,-2.2) rectangle (2.5,2.2);
      \draw[domain=0.001:2.2,samples=100,smooth,variable=\x]
        plot (\x, {pow(\x,1/3)*1.25});
      \draw[domain=-2.2:-0.001,samples=100,smooth,variable=\x]
        plot (\x, {pow(-\x,1/3)*1.25});
    \end{scope}
    \node[anchor=north] at (0.02,-0.10) {$O$};
    \node[anchor=south east] at (-0.05,1.25) {$1$};
    \node at (0.6,-2.85) {(D)};
  \end{scope}
\end{tikzpicture}
\end{center}
\nopagebreak[4]
\end{question}
\begin{question}
\gkti{7}定义在 $\mathbf R$ 上的函数 $f(x)$ 满足 $f(x)=\begin{cases}\log_2(4-x),&x\leqslant0,\\ f(x-1)-f(x-2),&x>0,\end{cases}$ 则 $f(3)$ 的值为（　）

（A）$-1$

（B）$-2$

（C）$1$

（D）$2$

\end{question}
\begin{question}
\gkti{8}设 $P$ 是 $\triangle ABC$ 所在平面内的一点，$\overrightarrow{BC}+\overrightarrow{BA}=2\overrightarrow{BP}$，则（　）

（A）$\overrightarrow{PA}+\overrightarrow{PB}=\mathbf 0$

（B）$\overrightarrow{PB}+\overrightarrow{PC}=\mathbf 0$

（C）$\overrightarrow{PC}+\overrightarrow{PA}=\mathbf 0$

（D）$\overrightarrow{PA}+\overrightarrow{PB}+\overrightarrow{PC}=\mathbf 0$

\end{question}
\begin{question}
\gkti{9}已知 $\alpha,\beta$ 表示两个不同的平面，$m$ 为平面 $\alpha$ 内的一条直线，则「$\alpha\perp\beta$」是「$m\perp\beta$」的（　）

（A）充分不必要条件

（B）必要不充分条件

（C）充要条件

（D）既不充分也不必要条件

\end{question}
\begin{question}
\gkti{10}设斜率为 $2$ 的直线 $l$ 过抛物线 $y^2=ax$（$a\neq0$）的焦点 $F$，且和 $y$ 轴交于点 $A$．若 $\triangle OAF$（$O$ 为坐标原点）的面积为 $4$，则抛物线方程为（　）

（A）$y^2=\pm4x$

（B）$y^2=\pm8x$

（C）$y^2=4x$

（D）$y^2=8x$

\end{question}
\begin{question}
\gkti{11}在区间 $\left[-\dfrac{\pi}{2},\dfrac{\pi}{2}\right]$ 上随机取一个数 $x$，$\cos x$ 的值介于 $0$ 到 $\dfrac12$ 之间的概率为（　）

（A）$\dfrac13$

（B）$\dfrac2\pi$

（C）$\dfrac12$

（D）$\dfrac23$

\end{question}
\begin{question}
\gkti{12}已知定义在 $\mathbf R$ 上的奇函数 $f(x)$ 满足 $f(x-4)=-f(x)$，且在区间 $[0,2]$ 上是增函数，则（　）

（A）$f(-25)<f(11)<f(80)$

（B）$f(80)<f(11)<f(-25)$

（C）$f(11)<f(80)<f(-25)$

（D）$f(-25)<f(80)<f(11)$

\end{question}
\gkpart{二、填空题}

\begin{question}
\gkti{13}在等差数列 $\{a_n\}$ 中，$a_3=7$，$a_5=a_2+6$，则 $a_6=$\gkblank{}．

\end{question}
\begin{question}
\gkti{14}若函数 $f(x)=a^x-x-a$（$a>0$ 且 $a\neq1$）有两个零点，则实数 $a$ 的取值范围是\gkblank{}．

\Needspace{0.4\textheight}
\end{question}
\begin{question}
\gkti{15}执行如图的程序框图，输出的 $T=$\gkblank{}．
\nopagebreak[4]
\begin{center}% 2009 年山东卷文科第 15 题：程序框图，输出 T 的值。
% 来源 PDF：普通高考/2009/2009山东文.pdf，SHA256 1807901a21da57e3d45382294189b0c8e795722b6879842319ae75428639fad6
% 原页：第 1 页（A3 三栏）第 3 栏上部，印刷页 547。
% 重绘依据：开始 → S=0,T=0,n=0 → 判断 T>S：是（右）输出 T 并结束；否（下）S=S+5 → n=n+2 → T=T+n，
%   再回到判断。
\begin{tikzpicture}[scale=0.9, semithick, >=Stealth,
  node distance=5mm, every node/.style={inner sep=3pt, font=\small}]
  \node (start) [draw, rounded corners=4pt, minimum width=1.9cm, minimum height=0.62cm] {开始};
  \node (init) [draw, minimum width=3.6cm, minimum height=0.62cm, below=of start] {$S=0,\ T=0,\ n=0$};
  \node (cond) [draw, diamond, aspect=2.2, minimum width=3.0cm, minimum height=0.95cm, below=of init] {$T>S$};
  \node (s) [draw, minimum width=2.2cm, minimum height=0.62cm, below=of cond] {$S=S+5$};
  \node (n) [draw, minimum width=2.2cm, minimum height=0.62cm, below=of s] {$n=n+2$};
  \node (t) [draw, minimum width=2.2cm, minimum height=0.62cm, below=of n] {$T=T+n$};
  \node (out) [draw, trapezium, trapezium left angle=70, trapezium right angle=110,
    minimum width=2.2cm, minimum height=0.62cm, right=1.6cm of cond] {输出 $T$};
  \node (end) [draw, rounded corners=4pt, minimum width=1.9cm, minimum height=0.62cm, below=of out] {结束};
  \draw[->] (start) -- (init);
  \draw[->] (init) -- (cond);
  \draw[->] (cond) -- node[right, pos=0.55] {否} (s);
  \draw[->] (s) -- (n);
  \draw[->] (n) -- (t);
  \draw[->] (t.west) -- ++(-1.3,0) |- (cond.west);
  \draw[->] (cond) -- node[above, pos=0.45] {是} (out);
  \draw[->] (out) -- (end);
\end{tikzpicture}
\end{center}
\nopagebreak[4]
\end{question}
\begin{question}
\gkti{16}某公司租赁甲、乙两种设备生产 $A,B$ 两类产品，甲种设备每天能生产 $A$ 类产品 $5$ 件和 $B$ 类产品 $10$ 件，乙种设备每天能生产 $A$ 类产品 $6$ 件和 $B$ 类产品 $20$ 件．已知设备甲每天的租赁费为 $200$ 元，设备乙每天的租赁费为 $300$ 元．现该公司至少要生产 $A$ 类产品 $50$ 件，$B$ 类产品 $140$ 件，所需租赁费最少为\gkblank{}元．

\end{question}
\gkpart{三、解答题}

\begin{question}
\gkti{17}已知函数 $f(x)=2\sin x\cos^2\dfrac x2+\cos x\sin\varphi-\sin x$（$0<\varphi<\pi$）在 $x=\pi$ 处取最小值．

（1）求 $\varphi$ 的值；

（2）在 $\triangle ABC$ 中，$a,b,c$ 分别是角 $A,B,C$ 的对边．已知 $a=1$，$b=\sqrt2$，$f(A)=\dfrac{\sqrt3}{2}$，求角 $C$．

\Needspace{0.4\textheight}
\end{question}
\begin{question}
\gkti{18}如图，在直四棱柱 $ABCD-A_1B_1C_1D_1$ 中，底面 $ABCD$ 为等腰梯形，$AB\parallel CD$，$AB=4$，$BC=CD=2$，$AA_1=2$，$E$、$E_1$ 分别是棱 $AD$、$AA_1$ 的中点．
\nopagebreak[4]
\begin{center}% 2009 年山东卷文科第 18 题：直四棱柱 ABCD-A1B1C1D1，底面为等腰梯形。
% 来源 PDF：普通高考/2009/2009山东文.pdf，SHA256 1807901a21da57e3d45382294189b0c8e795722b6879842319ae75428639fad6
% 原页：第 2 页（A3 三栏）第 1 栏上部，印刷页 548。
% 重绘依据：底面 ABCD 按斜投影画成等腰梯形（AB 为下底、CD 为上底），A 左下、B 右下、C 右上、D 左上；
%   A1B1C1D1 为对应上底；E 在 AD 上、E1 在 AA1 上、F 在 AB 上；
%   被遮挡的棱 A1D1、D1C1、DD1 与辅助线画虚线。
\begin{tikzpicture}[scale=0.95, semithick, every node/.style={inner sep=1.3pt}]
  \coordinate (A) at (0,0);
  \coordinate (B) at (4.0,0);
  \coordinate (C) at (2.9,1.15);
  \coordinate (D) at (1.1,1.15);
  \coordinate (A1) at (0,2.2);
  \coordinate (B1) at (4.0,2.2);
  \coordinate (C1) at (2.9,3.35);
  \coordinate (D1) at (1.1,3.35);
  \coordinate (E) at ($(A)!0.5!(D)$);
  \coordinate (E1) at ($(A)!0.5!(A1)$);
  \coordinate (F) at ($(A)!0.5!(B)$);
  % 底面
  \draw (A) -- (B) -- (C) -- (D) -- cycle;
  % 上底（部分被遮挡）
  \draw (A1) -- (B1) -- (C1);
  \draw[dashed] (C1) -- (D1) -- (A1);
  % 侧棱
  \draw (A) -- (A1);
  \draw (B) -- (B1);
  \draw (C) -- (C1);
  \draw[dashed] (D) -- (D1);
  % 点
  \fill (E) circle (1.1pt);
  \fill (E1) circle (1.1pt);
  \fill (F) circle (1.1pt);
  \node[anchor=north west] at (A) {$A$};
  \node[anchor=north east] at (B) {$B$};
  \node[anchor=west] at (C) {$C$};
  \node[anchor=south east] at (D) {$D$};
  \node[anchor=east] at (A1) {$A_1$};
  \node[anchor=south east] at (B1) {$B_1$};
  \node[anchor=south west] at (C1) {$C_1$};
  \node[anchor=east] at (D1) {$D_1$};
  \node[anchor=north west] at (E) {$E$};
  \node[anchor=east] at (E1) {$E_1$};
  \node[anchor=north] at ([yshift=-2pt]F) {$F$};
\end{tikzpicture}
\end{center}
\nopagebreak[4]
（1）设 $F$ 是棱 $AB$ 的中点，证明：直线 $EE_1\parallel$ 平面 $FCC_1$；

（2）证明：平面 $D_1AC\perp$ 平面 $BB_1C_1C$．

\Needspace{0.4\textheight}
\end{question}
\begin{question}
\gkti{19}一汽车厂生产 $A,B,C$ 三类轿车，每类轿车均有舒适型和标准型两种型号，某月的产量如下表（单位：辆）：
\nopagebreak[4]
\begin{center}% 2009 年山东卷文科第 19 题：三类轿车两型号的产量表。
% 来源 PDF：普通高考/2009/2009山东文.pdf，SHA256 1807901a21da57e3d45382294189b0c8e795722b6879842319ae75428639fad6
% 原页：第 2 页（A3 三栏）第 1 栏下部，印刷页 548。
% 重绘依据：表头为「车型」与轿车 A、B、C，两行为舒适型、标准型；A 列 100/300，B 列 150/450，
%   C 列舒适型为 z、标准型为 600；单位「辆」随题干给出。
\begin{tikzpicture}[scale=1.0, semithick, every node/.style={inner sep=3pt, font=\small}]
  \draw (0,0) rectangle (7.2,-1.6);
  \draw (0,-0.8) -- (7.2,-0.8);
  \draw (1.8,0) -- (1.8,-1.6);
  \draw (3.6,0) -- (3.6,-1.6);
  \draw (5.4,0) -- (5.4,-1.6);
  \node at (0.9,-0.4) {车型};
  \node at (2.7,-0.4) {轿车 $A$};
  \node at (4.5,-0.4) {轿车 $B$};
  \node at (6.3,-0.4) {轿车 $C$};
  \node at (0.9,-1.2) {舒适型};
  \node at (2.7,-1.2) {$100$};
  \node at (4.5,-1.2) {$150$};
  \node at (6.3,-1.2) {$z$};
  \draw (0,-2.4) -- (7.2,-2.4);
  \draw (0,0) -- (0,-2.4);
  \draw (1.8,0) -- (1.8,-2.4);
  \draw (3.6,0) -- (3.6,-2.4);
  \draw (5.4,0) -- (5.4,-2.4);
  \draw (7.2,0) -- (7.2,-2.4);
  \node at (0.9,-2.0) {标准型};
  \node at (2.7,-2.0) {$300$};
  \node at (4.5,-2.0) {$450$};
  \node at (6.3,-2.0) {$600$};
\end{tikzpicture}
\end{center}
\nopagebreak[4]
按类型分层抽样的方法在这个月生产的轿车中抽取 $50$ 辆，其中有 $A$ 类轿车 $10$ 辆．

（1）求 $z$ 的值；

（2）用分层抽样的方法在 $C$ 类轿车中抽取一个容量为 $5$ 的样本．将该样本看成一个总体，从中任取 $2$ 辆，求至少有 $1$ 辆舒适型轿车的概率；

（3）用随机抽样的方法从 $B$ 类舒适型轿车中抽取 $8$ 辆，经检测它们的得分如下：$9.4$，$8.6$，$9.2$，$9.6$，$8.7$，$9.3$，$9.0$，$8.2$．把这 $8$ 辆轿车的得分看成一个总体，从中任取一个数，求该数与样本平均数之差的绝对值不超过 $0.5$ 的概率．

\end{question}
\begin{question}
\gkti{20}等比数列 $\{a_n\}$ 的前 $n$ 项和为 $S_n$．已知对任意的 $n\in\mathbf N^*$，点 $(n,S_n)$ 均在函数 $y=b^x+r$（$b>0$ 且 $b\neq1$，$b,r$ 为常数）的图象上．

（1）求 $r$ 的值；

（2）当 $b=2$ 时，记 $b_n=\dfrac{n+1}{4a_n}$（$n\in\mathbf N^*$），求数列 $\{b_n\}$ 的前 $n$ 项和 $T_n$．

\end{question}
\begin{question}
\gkti{21}已知函数 $f(x)=\dfrac13ax^3+bx^2+x+3$，其中 $a\neq0$．

（1）当 $a,b$ 满足什么条件时，$f(x)$ 取得极值？

（2）已知 $a>0$，且 $f(x)$ 在区间 $(0,1]$ 上单调递增，试用 $a$ 表示 $b$ 的取值范围．

\end{question}
\begin{question}
\gkti{22}设 $m\in\mathbf R$，在平面直角坐标系中，已知向量 $\boldsymbol a=(mx,y+1)$，向量 $\boldsymbol b=(x,y-1)$，$\boldsymbol a\perp\boldsymbol b$，动点 $M(x,y)$ 的轨迹为 $E$．

（1）求轨迹 $E$ 的方程，并说明该方程所表示曲线的形状；

（2）已知 $m=\dfrac14$，证明：存在圆心在原点的圆，使得该圆的任意一条切线与轨迹 $E$ 恒有两个交点 $A,B$，且 $OA\perp OB$（$O$ 为坐标原点），并求出该圆的方程；

（3）已知 $m=\dfrac14$，设直线 $l$ 与圆 $C:x^2+y^2=R^2$（$1<R<2$）相切于 $A_1$，且 $l$ 与轨迹 $E$ 只有一个公共点 $B_1$，当 $R$ 为何值时，$|A_1B_1|$ 取得最大值？并求最大值．

\end{question}